\documentclass[11pt,a4paper]{article}
\usepackage[margin=2.5cm]{geometry}
\usepackage{amsmath,amssymb}
\usepackage{booktabs}
\usepackage{microtype}
\usepackage{fancyvrb}
\usepackage[hidelinks]{hyperref}
\DefineVerbatimEnvironment{board}{Verbatim}{fontsize=\normalsize,xleftmargin=2em}

\title{Fortresses in Wolf and Lamb}
\author{}
\date{26 September 2026}

\begin{document}
\maketitle

\section*{Summary}

Both sides have a way of never losing, but they are of different kinds.

\begin{itemize}
\item \textbf{The lambs have true fortresses: shapes that cannot be broken whatever
the wolves do.} None exists with 4 or 5 lambs; with 6 lambs there are exactly four
(up to symmetry), with 7 lambs 141, with 8 lambs 1{,}179. They all follow one
rule, the \emph{distance-two rule}: a lamb can only be captured from a square two
steps away in its row or column, so a lamb whose distance-two squares are all
either off the board or occupied by lambs is untouchable. The catch is that
lambs must move every turn, so a fortress also needs spare moves that do not open
a hole. The 4-lamb square that satisfies the rule perfectly is nevertheless
always lost, by zugzwang.
\item \textbf{The wolves have no fortress shape.} Every wolf placement can be
trapped for some lamb configuration. What the wolves have instead is a
\emph{principle}: stay apart and off the edges. With 9 lambs left, wolves that
are not adjacent to each other and not on the edge lose in 0.9\,\% of positions;
wolves huddled on the edge lose in half of them. The reason is counting: a corner
wolf is trapped by 2 lambs, an edge wolf by 3, a central wolf by 4, and three
separate nets cost more lambs than the flock can spare.
\item In the drawn starting line the two meet: the wolves eat down to a few lambs
that have settled into a distance-two shape, the wolves are spread out, and
nobody can make progress.
\end{itemize}

All statements below are read off the exact solution tables: for a shape, ``never
lost'' means that over every placement of the three wolves and both sides to
move, no position is a win for the other side.

\section*{The distance-two rule}

A wolf captures by jumping over one empty square onto a lamb. So a lamb on square
$s$ can be captured along a line only if the neighbour of $s$ on that line is
empty and the square two steps away on the same side is on the board and reachable
by a wolf. This gives:

\begin{quote}
\textbf{A lamb is immune along a line if, on each side, the square two steps away
is off the board or holds a lamb, or the adjacent square is occupied.}
\end{quote}

The second condition is fragile (the occupying piece may move), the first is
stable as long as the protecting lamb stays. Squares two steps apart in a row or
column have the same parity of row and of column, so the board splits into four
\emph{lattices}: the 9 squares with both coordinates odd, the 4 squares with both
even, and two lattices of 6 with mixed parity. A set of lambs that fills a whole
lattice is immune while it stands still:

\begin{board}
   even-even (4)        odd-even (6)         even-odd (6)
   . . . . .            . L . L .            . . . . .
   . L . L .            . . . . .            L . L . L
   . . . . .            . L . L .            . . . . .
   . L . L .            . . . . .            L . L . L
   . . . . .            . L . L .            . . . . .
\end{board}

\section*{Why the 4-lamb square is not a fortress}

The four-square lattice is perfectly immune, yet the tables show it is lost in
every position where the wolves are close. Lambs may not pass. Whatever lamb
moves, it leaves its lattice square, loses the protection of its distance-two
partners, and the lattice has no spare lamb to fill the gap; the wolves then win
in a few moves. Example (wolf to move, wolves win in every line; also lost with
the lambs to move, all 16 lamb moves losing):

\begin{board}
   W . . . .
   . L . L .
   . . W . .
   . L . L .
   . . . . W
\end{board}

With 4 lambs no shape at all is a fortress: one capture ends the game, and the
best shape is still lost in 81\,\% of positions. With 5 lambs the best shape is
lost in 36\,\%.

\section*{The 6-lamb fortresses}

Six lambs is the smallest number that admits shapes which are never lost. There
are exactly four up to symmetry:

\begin{board}
   A               B               C               D
   . L . L .       . L . L .       . . . . .       . . . . .
   . L . L .       . . . . .       L . L . L       . L . L .
   . . . . .       . L . L .       . . . . .       . L . L .
   . L . L .       . L . L .       L . L . L       . L . L .
   . . . . .       . . . . .       . . . . .       . . . . .
\end{board}

C is the full even-odd lattice (6 lambs) and is drawn in all 524 positions. A, B
and D are \emph{double rails}: two columns one square in from the edges, three
lambs each, with two of the three lambs stacked. Each rail protects the other at
distance two horizontally; inside a rail the stacked pair protects itself by
occupancy and the third lamb by distance two. A and B are drawn in 1{,}012 of
1{,}014 positions and won by the lambs in the other 2; nothing is ever lost.

What makes these six-lamb shapes work where the square failed is \emph{tempo}.
In shape A with the wolves probing as below, the lambs to move have 17 legal
moves, 13 of which keep the draw:

\begin{board}
   . L . L .
   . L . L .
   W . W . .        Lamb to move: draw.
   . L . L .        Drawing moves include (1,2)-(1,1), (1,4)-(1,3),
   . . . W .        (2,2)-(2,1), (2,4)-(2,3), (4,2)-(4,1), (2,2)-(3,2), ...
\end{board}

A lamb steps sideways off its rail and back again; while it is away, its
neighbours still cover the squares two steps from it, so the wolves gain nothing.
The four losing moves are the ones that put a lamb two steps from a wolf with an
empty square between them, such as (4,4)-(3,4) next to the wolf on (3,3) or
(4,2)-(5,2) in line with the wolf on (5,4).

Because the rule is local, the fortresses grow: with 7 lambs the extra lamb
usually sits in a corner or on an edge square that the rails already protect, and
with 8 lambs the most robust shapes put a full rail on the top edge:

\begin{board}
   7 lambs                     8 lambs
   . . . L L                   L L L . L
   . L . L .                   . L . L .
   . . . . .                   . . . . .
   . L . L .                   . L . L .
   L . . . .                   . . . . .
\end{board}

Shapes that look just as tidy but break the rule are lost. The best five-lamb
shape, a row of three with two supporters two rows below, is still lost in
36\,\% of positions; the best four-lamb shapes in 81 to 87\,\%.

\section*{The wolves: no shape, but a principle}

For every placement of the three wolves, some lamb configurations trap them,
even with only 7 lambs. So the wolves have no fortress in the strict sense. The
tables do, however, show a strong pattern in \emph{which} placements survive.
Over all positions with 9 lambs:

\begin{center}
\begin{tabular}{@{}llrrr@{}}
\toprule
adjacent wolf pairs & wolves on the edge & wolf wins & draws & lamb wins \\
\midrule
0 & 0 & 54.3\,\% & 44.8\,\% & 0.9\,\% \\
0 & 1 & 36.5\,\% & 59.2\,\% & 4.3\,\% \\
0 & 2 & 18.6\,\% & 68.1\,\% & 13.3\,\% \\
0 & 3 & 7.7\,\% & 62.1\,\% & 30.2\,\% \\
1 & 3 & 9.5\,\% & 50.6\,\% & 39.9\,\% \\
2 & 3 & 12.8\,\% & 36.2\,\% & 50.9\,\% \\
\bottomrule
\end{tabular}
\end{center}

Two effects, both large. Every wolf on the edge costs the wolves about a third
of their winning chances and multiplies their losing chances. And wolves that
touch each other are easier to trap than wolves that stay apart: with all three
on the edge, going from no adjacent pair to two adjacent pairs raises the lamb
win rate from 30\,\% to 51\,\%. A huddle can be closed with one net; three
separated wolves need three nets, and the flock cannot afford them. The safest
placements in the middle game are spread as far apart as the board allows:

\begin{board}
   7 lambs, 87% draws     8 lambs, 89% draws     9 lambs, 85% draws
   W . . . W              . . . . W              . . . . .
   . . . . .              . . . . .              . . . . .
   . . . . .              . . W . .              W . W . W
   . . . . .              . . . . .              . . . . .
   W . . . .              W . . . .              . . . . .
\end{board}

The count behind the principle: to immobilise a wolf the lambs must occupy all its
empty neighbours \emph{and} keep the square beyond each neighbour from giving a
jump. In a corner that takes 2 lambs, on an edge 3, in the open 4, and the
lambs also need tempo moves of their own. With three wolves far apart the bill
exceeds what 7 to 9 lambs can pay, and the game is a draw.

\section*{How the two fortresses meet in the main line}

From the start the only drawing wolf move is the central capture. The wolves
then win material through double threats while the lambs give up lambs that
cannot be saved and keep the ones that can. Best play by both sides reaches
positions with a handful of lambs in a distance-two shape and three wolves spread
across the board: the lambs are immune, the wolves are untrappable, and the
100-turn rule ends the game. One such ending, reached with the tables playing
both sides:

\begin{board}
   W . W . .
   L . L . .
   L . . L .        Wolf to move, draw. Every lamb is covered by
   W . . . .        a partner or an edge; every wolf has room.
   . . . . .
\end{board}

\section*{Practical advice}

\begin{itemize}
\item \textbf{Lambs:} think in distance two. A lamb is safe if the squares two
away in its row and column are lambs or off the board. Build rails a square in
from the edge, keep pairs stacked, and always keep a harmless waiting move in
reserve; if every move would open a hole, the shape is a zugzwang, not a
fortress. Never trade the last spare lamb: six is the minimum that can hold.
\item \textbf{Wolves:} never let two wolves touch, keep at least one wolf in the
centre, and stay off the edges unless a capture forces it. A wolf that is
tempted to the edge by a free lamb has often walked into the only net the lambs
can afford.
\end{itemize}

\section*{How this was computed}

\texttt{solver/fortress.js lamb $k$} groups every position with $k$ lambs by the
lamb shape (up to the eight symmetries of the board) and counts wolf wins, lamb
wins and draws over all wolf placements and both sides to move; shapes with zero
wolf wins are the fortresses. \texttt{solver/fortress.js wolf $k$} does the same
by wolf placement, and \texttt{solver/wolfshape.js $k$} tabulates the value by
the wolves' formation. The immunity of individual lambs was checked directly
from the capture rule.

\end{document}
